Com $g(h)=g_0 (R/(R+h))^{2}$, o período do pêndulo varia como
\[
T(h)=2\pi\sqrt{\frac{L}{g(h)}} = T_0\frac{R+h}{R}.
\]
Assim, o período cresce praticamente em proporção linear com a altura:
\[
T(h) \approx T_0\left(1 + \frac{h}{R}\right).
\]